\section{Fremstilling}

\subsection{Fumlebræt-modeller}
	\begin{figure}[H]
		\centering
	    \includegraphics[width=13cm]{figures/stock.jpg}
		\caption{}
		\label{fig:}
	\end{figure}
\subsection{PCB - fremstilling}\label{subs:pcbfremstilling}
	\begin{figure}[H]
		\centering
	    \includegraphics[width=13cm]{figures/stock.jpg}
		\caption{Et billede af et PCB vi designede i Livewire}
		\label{fig:PCBPrint}
	\end{figure}
	\begin{figure}[H]
		\centering
	    \includegraphics[width=13cm]{figures/stock.jpg}
		\caption{Et billede af et PCB fremstillet fysisk}
		\label{fig:PCBfysisk}
	\end{figure}
\subsection{Endelige prototype}\label{subs:endeligProto}
\begin{figure}[H]
	\centering
    \includegraphics[width=13cm]{figures/stock.jpg}
	\caption{Billede af endelig produkt}
	\label{fig:endeligPrototype}
\end{figure}
\todo{Beskriv hvordan PCB-baseret produkt ikke fungrede, har vi nogen mulige forklaringer?}
punkt 1 PCB'et havde muligvist ikke ordenlige forbindelser
Dette har vi testet ved at bruge et multimeter se om der var gode forbindelser mellem alle punkter på PCB'et
punkt 2 Der er mulighed for at da PCB'et manglede en ledning eller 2 da det blev lavet på computeren.
Dette er noget som vi har prøvet at undgå ved at lave så få manuelle ændringer i PCB-wizard så der ikke blev overset nogle forbindelser. 
punkt 3 Der var komponeter som ikke fungerede ordenligt.
Det er ikke særlig ofte at komponeter ikke fungere ordenligt når vi bruger dem derfor kan det væod at der er en komponent der er blev sat forkert i eller at der er blevet brugt et forkert komponent hvilket også er en del af dette punkt.

